\subsection{Approche variationnelle}
\subsubsection{L'équation scalaire}
Cette section fait suite à la section \ref{RS_scalaire}. Nous allons ici suivre l'analyse mathématique proposée par Showalter et Ting  \cite{Showalter_Ting} pour étudier l'équation scalaire \eqref{eq_RS_scalaire}. Comme mentionné dans l'introduction, nous allons considérer un domaine borné de $\bbR^d$. Ce choix est motivé par les problèmes physiques qui se posent dans des systèmes réels, simulés sur des domaines finis. Pour simplifier l'étude, on prendra des conditions de Dirichlet homogènes. 
Le cadre fonctionnel approprié est donc $H^1_0$. 
On considère un ouvert borné $\Omega \subset \bbR^d$.
Notre inconnue $u$ dépend à la fois du temps et de l'espace, pour trouver une formulation faible de notre 
problème \eqref{eq_RS_scalaire} on doit donc gérer la dépendance en temps. Pour ça, on va \og regarder \fg{}
notre inconnue $\bbR \times \bbR^d \ni (t,x) \mapsto u(t,x) \in \bbR$ comme une fonction 
$\bbR \ni t \mapsto \tilde{u}(t)$ à valeurs dans un espace vectoriel de fonctions définies sur $\bbR$, c'est-à-dire 
$\bbR^d \ni x \mapsto \tilde{u}(t)(x) \in \bbR$. Ici, on a un problème d'ordre $2$ en espace avec des conditions 
de Dirichlet homogènes, on va donc chercher une solution $u : \bbR \rightarrow \Hone$.
Si on se donne une donnée initiale $u_0$, on peut reformuler \eqref{eq_RS_scalaire} en : trouver $u \in \scrC^1(\bbR, \Hone)$ tel que pour tout $t \in \bbR$
\begin{equation}\label{pb_RS}
\begin{cases}
    u'(t) - \nu\Delta u(t) - \alpha\Delta \big(u'(t)\big) = 0  \\
    u(0) = u_0
\end{cases}
\end{equation}
On définit deux opérateurs $M$ et $L$ par
\begin{eqnarray*}
    M := I - \alpha \Delta \ \ \ \ \ &\text{ et }& \ \ \ \ \ L :=-\nu\Delta
\end{eqnarray*}
de sorte à reformuler le problème \eqref{pb_RS} en : trouver $u \in \scrC^1(\bbR, \Hone)$ tel que pour tout $t \in \bbR$
\begin{equation}\label{RS_forme_operateur}
\begin{cases}
    Mu'(t) + Lu(t)=0  \\
    u(0) = u_0
\end{cases}
\end{equation} 
Nous allons à présent étudier les opérateurs $M$ et $L$ pour montrer l'existence et l'unicité d'une solution à notre problème.

\subsubsection*{Extension de \texorpdfstring{$M$}{M} et \texorpdfstring{$L$}{L}.}
A priori, nos opérateurs sont définis de $\scrC^2(\Omega)$ dans $\scrC^0(\Omega)$. Pour étendre leurs domaines de définition nous allons utiliser le théorème suivant, du à Lax et Milgram \cite[p.~171]{theory_PDE}.

\begin{theoreme}\label{thm_lax}: \\
$\ulcorner$Soient $(H, \langle \cdot, \cdot \rangle)$ et $(H_0, \langle \cdot, \cdot \rangle_0)$ deux espaces de Hilbert tels que $H$ est un sous-espace dense de $H_0$ et 
\begin{equation}\label{ineq_thm_lax}
    \exists k >0,\ \forall x \in H, \ \norm{x} \geq k \norm{x}_0
\end{equation}
Soit $B$ une forme bilinéaire continue et coercive sur $H$. Alors il existe deux opérateurs linéaires $E$ et $E^*$, dont les domaines de définitions $D$ et $D^*$ sont denses dans $H$, tels que
\begin{enumerate}
    \item On a
    \begin{eqnarray*}
        \forall x \in D,\ \forall y \in H, \ B(x,y) &=& \ps{Ex}{y}_0 \\
        \forall y \in D^*,\ \forall x \in H, \ B(x,y) &=& \ps{x}{E^* y}_0
    \end{eqnarray*}
    \item L'image de $E$ est $H_0$, de même pour $E^*$.
    \item $E$ et $E^*$ admettent des inverses $S$ et $S^*$ qui sont continus de $H_0$ dans $H$.
    \begin{eqnarray*}
        \exists k_1>0,\ \forall z \in H_0, \ \norm{Sz} &\leq& k_1 \norm{z}_0 \\
        \exists k_1^*>0,\ \forall z \in H_0, \ \norm{S^*z} &\leq& k_1^* \norm{z}_0
    \end{eqnarray*}
    Et donc, $S$ et $S^*$ sont continus de $H_0$ dans $H_0$
    \begin{eqnarray*}
        \exists k_2>0,\ \forall z \in H_0, \ \norm{Sz}_0 &\leq& k_2 \norm{z}_0 \\
        \exists k_2^*>0,\ \forall z \in H_0, \ \norm{S^*z}_0 &\leq& k_2^* \norm{z}_0
    \end{eqnarray*}
\end{enumerate}
\end{theoreme}

Un démonstration du théorème est proposée en Annexe \ref{AnnexeA}. Avant de commencer, quelques remarques sur ce théorème.

\begin{remark}
    \begin{enumerate}
        \item Le cadre : on considère deux espaces de Hilbert $H$ et $H_0$ dont l'un est \og plus fin \fg{} que l'autre, dans le sens où il est inclus dans l'autre et sa norme est \og plus restrictive \fg{}. On considère ensuite une forme bilinéaire $B$ continue et coercive sur $H$, donc la forme bilinéaire se comporte très bien sur $H$.
        \item L'idée générale : de la même manière que le théorème de représentation de Riesz, on va chercher à \og représenter \fg{} la forme bilinéaire $B$ par un produit scalaire. Cependant, l'idée va un peu plus loin puisqu'on veut la représenter sur un espace \og plus grand \fg{} et par un opérateur linéaire.
        \item Le but : passer d'un problème variationnel à une équation de la forme $Eu = f$ pour utiliser la théorie des semi-groupes ou encore la théorie spectrale.
        \item Le premier point du théorème permet donc de \og représenter \fg{} la forme bilinéaire $B$ à l'aide d'un opérateur linéaire $E$ et d'un produit scalaire tout en \og l'étendant \fg{} dans le sens où le produit scalaire est celui sur $H_0$ qui est le \og plus grand \fg{} espace. On se contente ici de parler de $E$, les arguments pour $E^*$ étant les mêmes. Comme la forme bilinéaire $B$ n'est pas symétrique on a $E^* \neq E$ et $E^*$ est l'adjoint de $E$.
        \item Le deuxième point du théorème nous indique que l'opérateur $E$ est bijectif, ce qui signifie que pour tout second membre $f \in H_0$ on a l'existence d'une unique solution à l'équation $Eu=f$.
        \item Le troisième point du théorème nous assure l'existence d'un inverse $S$ de $E$, ainsi l'équation $Eu=f$ admet pour solution $u=Sf$. La continuité de l'opérateur $S$ nous garantit le bon comportement de la solution par rapport à la donnée initiale. De plus, la continuité de $S : H_0 \rightarrow H$ nous permet de gagner en \og régularité \fg{}.
    \end{enumerate} 
\end{remark}

Dans notre cas, on considère les espaces de Hilbert $H=H^1_0(\Omega)$ et $H_0=L^2(\Omega)$. On utilisera les inclusions suivantes
\begin{equation*}
    \scrD \ \subset \ \Hone \ \subset \ L^2(\Omega)
\end{equation*}
et comme $\overline{\scrD}^{L^2(\Omega)}$ = $L^2(\Omega)$, on a bien 
\begin{equation*}
    \overline{\Hone}^{L^2(\Omega)} = L^2(\Omega)
\end{equation*}
De plus, on a 
\begin{equation*}
    \forall \varphi \in \Hone, \ \|\varphi\|_{L^2} \leq \|\varphi\|_{H^1}
\end{equation*}
On va maintenant chercher à établir une formulation faible de notre équation du problème \eqref{pb_RS}. Soit donc $v \in \Hone$.
\begin{eqnarray*}
    \int_{\Omega}Mu'(t)v + \int_{\Omega} Lu(t)v =0 &\Longrightarrow&
    \int_{\Omega} u'(t)v - \alpha \int_{\Omega} \Delta u'(t)v - \nu \int_{\Omega}
    \Delta u(t)v = 0 \\
    &\overset{\text{Green}}{\Longrightarrow}& \int_{\Omega} u'(t)v + \alpha \int_{\Omega} \nabla u'(t)\cdot\nabla v + \nu \int_{\Omega}
    \nabla u(t)\cdot \nabla v = 0 \\
    &\Longrightarrow& \int_{\Omega} \big(u'(t)v+\alpha\nabla u'(t)\cdot\nabla v\big)
    + \nu\int_{\Omega} \nabla u(t)\cdot \nabla v = 0
\end{eqnarray*}
On définit ainsi deux formes bilinéaires sur $\Hone$
\begin{eqnarray*}
    b_M(\varphi, \psi) = \int_{\Omega} \big(\varphi\psi+\alpha\nabla \varphi\cdot\nabla \psi\big) \ \ \ \ \ &\text{ et }& \ \ \ \ \
    b_L(\varphi, \psi) = \nu\int_{\Omega} \nabla \varphi\cdot \nabla \psi
\end{eqnarray*}

\begin{remark}
    Si $\varphi$ et $\psi$ sont suffisamment régulières (par exemple $\varphi, \psi \in \scrD$), alors les implications précédentes sont des équivalences et donc on a
    \begin{eqnarray*}
        b_M(\varphi,\psi) = \ps{M\varphi}{\psi}_{L^2} \ \ \ \ \ &\text{ et }& \ \ \ \ \ b_L(\varphi,\psi) = \ps{L\varphi}{\psi}_{L^2}
    \end{eqnarray*}
\end{remark}

On obtient ainsi le problème variationnel associé au problème continu \eqref{RS_forme_operateur} : pour tout $t \in \bbR$, trouver $u(t) \in \Hone$ tel que
\begin{equation*}
    \forall v \in \Hone, \ \frac{\mathrm{d}}{\mathrm{d}t} b_M\big(u(t),v\big) + b_L\big(u(t),v\big) = 0
\end{equation*}
Nous cherchons maintenant à appliquer notre théorème. Montrons que $b_M$ et $b_L$ sont continues et coercives sur $H^1_0(\Omega)$.
\begin{itemize}
    \item Soient $\varphi, \psi \in \Hone$.
    \begin{align*}
        \abs{ b_M(\varphi,\psi)} &= \abs{\int_{\Omega}\varphi \psi + \alpha \int_{\Omega}\nabla\varphi\cdot\nabla\psi} \\
        & \overset{\text{C-S}}\leq \norm{\varphi}_{L^2}\norm{\psi}_{L^2} + \alpha \norm{\nabla\varphi}_{L^2}\norm{\nabla\psi}_{L^2}  \\
        &\leq (1+\alpha)\norm{\varphi}_{H^1}\norm{\psi}_{H^1}
    \end{align*}
    Donc $b_M$ est continue sur $H^1_0(\Omega)$.
    \item Soit $\varphi \in \Hone$.
    \begin{equation*}
        b_{M}(\varphi,\varphi) = \int_{\Omega} \varphi^2 + \alpha \int_{\Omega}\abs{\nabla \varphi}^2 \geq \min(1,\alpha) \norm{\varphi}^2_{H^1}
    \end{equation*}
    Donc $b_M$ est coercive sur $H^1_0(\Omega)$
    \item Soient $\varphi, \psi \in \Hone$.
    \begin{align*}
        \abs{ b_L(\varphi,\psi)} &= \abs{\nu \int_{\Omega}\nabla\varphi\cdot\nabla\psi}\\
        & \overset{\text{C-S}}\leq \nu \norm{\nabla\varphi}_{L^2}\norm{\nabla\psi}_{L^2} \\
        &\leq \nu \norm{\varphi}_{H^1}\norm{\psi}_{H^1}
    \end{align*}
    Donc $b_L$ est continue sur $H^1_0(\Omega)$.
    \item Soit $\varphi \in \Hone$.
    \begin{equation*}
        b_{L}(\varphi,\varphi) =\nu \norm{\nabla \varphi}^2_{L^2}
        =\nu \abs{\varphi}^2_{H^1} 
        \geq C_1 \nu \norm{\varphi}^2_{H^1}
    \end{equation*}
    où $|\cdot|_{H^1}$ désigne la norme $\Hone$, et on utilise le fait que sur $\Hone$, $\|\cdot\|_{H^1}$ et $|\cdot|_{H^1}$ sont équivalentes, donc $C_1 \in \bbR$ est un constante. Donc $b_L$ est coercive sur $H^1_0(\Omega)$.
\end{itemize}
On peut appliquer le théorème \ref{thm_lax}, ainsi il existe deux opérateurs $M_0$ et $L_0$, définis respectivement sur $D(M_0)$ et $D(L_0)$, qui sont denses dans $H^1_0(\Omega)$, tels que
\begin{align*}
    &\forall \varphi \in D(M_0),\ \forall \psi \in H^1_0(\Omega), \ b_M(\varphi, \psi) = \ps{M_0\varphi}{\psi}_{L^2} \\
    &\forall \varphi \in D(L_0),\ \forall \psi \in H^1_0(\Omega), \ b_M(\varphi, \psi) = \ps{L_0\varphi}{\psi}_{L^2}
\end{align*}
En particulier, on a pour tout $\varphi \in \scrC^2(\Omega)$ et pour tout $\psi \in H^1_0(\Omega)$ que
$$ \ps{M\varphi}{\psi}_{L^2} - \ps{M_0\varphi}{\psi}_{L^2} = 0 \ \Longrightarrow \ \ps{(M-M_0)\varphi}{\psi}_{L^2} $$
En prenant $\psi = (M-M_0)\varphi$ on obtient
\begin{equation*}
    \norm{(M-M_0)\varphi}^2_{L^2} = 0 \ \Longrightarrow \ (M-M_0)\varphi = 0
    \ \Longrightarrow \ M\varphi = M_0 \varphi
\end{equation*}
donc $M_0=M$ sur $\scrC^2(\Omega)$ ce qui montre bien que $M_0$ est une extension de $M$. Un raisonnement identique nous donne $L_0 = L$ sur $\scrC^2(\Omega)$ et donc $L_0$ est bien une extension de $L$. Le théorème nous assure également de l'existence de $M^{-1}_0$ et du fait que cet inverse est continu de $L^2(\Omega)$ dans $H^1_0(\Omega)$. Finalement, \eqref{RS_forme_operateur} se réécrit comme le problème suivant
\begin{equation}\label{edo_ML}
\begin{cases}
    u'(t) = -M^{-1}_0 L_0u(t)  \\
    u(0) = u_0
\end{cases}
\end{equation}

\subsubsection*{Extension de \texorpdfstring{$-M^{-1}_0 L_0$}{M-1\_0 L\_0}.}
Nous voulons désormais étendre $-M^{-1}_0 L_0$ pour que \eqref{RS_forme_operateur} soit bien défini dans $\Hone$. L'application du théorème nous a donne deux opérateurs étendus qui sont des bijections
\begin{eqnarray*}
    M_0 : D(M_0) \longrightarrow L^2(\Omega) \ \ \ \ \ &\text{ et }& \ \ \ \ \ L_0 : D(L_0) \longrightarrow L^2(\Omega)
\end{eqnarray*}
Cela justifie que  $-M^{-1}_0L_0 : D(L_0)\longrightarrow D(M_0)$ est une bijection. Pour étendre notre opérateur nous utilisons alors le résultat suivant vu dans le cours d'Analyse fonctionnelle.

\begin{theoreme}\label{thm_prolongement}
    Soient E et F deux espaces vectoriels normés. On suppose que F est complet et on considère un sous-espace dense $G \subset E$. Soit $f : G \longrightarrow F$ une application uniformément continue. Alors il existe une unique application $\mathscr{F} : E \longrightarrow F$ uniformément continue telle que $\mathscr{F}_{|G} = f$.
\end{theoreme}

On considère l'application $-M^{-1}_0L_0 : \scrD\longrightarrow \Hone$ qui aura le rôle de $f$ du théorème précédent. $\Hone$ est complet puisque c'est un Hilbert et $\scrD$ est dense dans $\Hone$ par définition. Il ne nous reste plus qu'à montrer que notre opérateur linéaire est uniformément continu, ce qui est équivalent à montrer qu'il est continu par linéarité.
\\
Soit $\varphi \in \scrD$. Alors 
\begin{eqnarray*}
    u = -\operateurs \varphi &\iff& M_0u = -L_0 \varphi
\end{eqnarray*}
En utilisant les inégalités de continuité sur $L_0$ et de coercivité sur $M_0$, on obtient
\begin{eqnarray*}
    \min(1,\alpha) \norm{u}^2_{H^1} \leq \ps{M_0u}{u}_{L^2} = \ps{-L_0\varphi}{u}_{L^2} \leq \nu \norm{\varphi}_{H^1}\norm{u}_{H^1}
\end{eqnarray*}
Ce qui revient à écrire 
\begin{eqnarray*}
    \min(1,\alpha) \norm{\operateurs \varphi}^2_{H^1} &\leq& \nu \norm{\varphi}_{H^1}\norm{\operateurs \varphi}_{H^1}
\end{eqnarray*}
et donc 
\begin{eqnarray*}
    \norm{\operateurs \varphi}_{H^1} \leq \frac{\nu}{\min(1,\alpha)}\norm{\varphi}_{H^1}
\end{eqnarray*}
Ainsi, l'opérateur $-\operateurs$ est continue sur $\Hone$. Par le théorème \eqref{thm_prolongement} il existe un unique opérateur $B : \Hone \longrightarrow \Hone$ qui prolonge $-\operateurs$.
\\
Pour l'instant rien n'indique que $B$ soit bijectif. Un candidat naturel pour son inverse est l'extension de $-L^{-1}_0M_0 : D(M_0) \longrightarrow D(L_0)$. En appliquant le même raisonnement qui vient de suivre on montre que l'on peut étendre $-L^{-1}_0M_0$ et que cette extension est précisément $B^{-1}$, de plus elle est continue sur $\Hone$. 
\\
Notre équation \eqref{edo_ML} devient donc
\begin{eqnarray}\label{edo_ML_B}
\begin{cases}
    u'(t) = Bu(t)  \\
    u(0) = u_0
\end{cases}
\end{eqnarray}

\subsubsection*{Existence et unicité.} 
Dans le cas d'un système différentiel linéaire de la forme
\begin{equation*}
    \left\{\begin{array}{c}
        u' + Au = 0 \\
        u(0) = u_0
    \end{array}\right.
\end{equation*}
Nous savons que la solution est de la forme $u = e^{-tA} u_0$. Pour résoudre \eqref{edo_ML_B} nous allons définir $\exp(tB)$. Tout d'abord, $B \in \mathscr{L}(\Hone)$. Et $\mathscr{L}(\Hone)$ est un espace vectoriel, donc pour tout $t \in \bbR$, $t B \in \mathscr{L}(\Hone)$. De plus, comme $\Hone$ est complet, $\mathscr{L}(\Hone)$ l'est aussi et on peut donc définir l'exponentielle de la façon suivante
\begin{eqnarray*}
\forall t \in \mathbb{R},\  E(t) := \exp(tB) = \sum_{k=0}^{+\infty} \frac{(tB)^k}{k!}
\end{eqnarray*}
À l'aide de la norme opérateur, on obtient
\begin{equation*}
    \forall t\in \mathbb{R}, \ \vertiii{E(t)}_{1} \leq \sum_{k=0}^{+\infty}\frac{\vertiii{(tB)^k}_{1}}{k!} \leq \sum_{k=0}^{+\infty}\frac{|t|^k\vertiii{B}_{1}^k}{k!} \leq \exp\big(\abs{t}\vertiii{B}_1\big) < +\infty
\end{equation*}
où $\vertiii{\cdot}_1$ désigne la norme d'opérateur sur $H^1$. Donc $\{ E(t) \ | \ t \in \mathbb{R} \} $ est un ensemble d'opérateurs continus de $\Hone$.
\begin{remark}\label{remarque_operateur}
    \begin{enumerate}
        \item \label{E_groupe_abelien} $\{ E(t) \ | \ t \in \mathbb{R} \}$ est un groupe abélien. 
        \item Chaque $E(t)$ est un opérateur bornée de $\Hone$ et la dépendance en $t$ est continue pour la topologie des opérateurs. En effet, soit 
        $t_0 \in \bbR$, on a par continuité de $\exp$ et par continuité de $B$
        \begin{eqnarray*}
            \vertiii{E(t)-E(t_0)}_1 &=& \vertiii{\ex^{tB}-\ex^{t_0B}}_1 \\
            &=& \vertiii{\ex^{t_0B}\big(\ex^{(t-t_0)B}-I\big)}_1 \\
            &\leq& \vertiii{\ex^{t_0B}}_1\vertiii{\ex^{(t-t_0)B}-I}_1 \\
            &=& \vertiii{\ex^{t_0B}}_1 \vertiii{\sum_{k=1}^{+\infty}\frac{\big((t-t_0)B\big)^k}{k!}} \\
            &\leq& \vertiii{\ex^{t_0B}}_1 (\exp{(|t-t_0|\vertiii{B}_1)}-1)\underset{t \to t_0}{\longrightarrow} 0
        \end{eqnarray*}
        \item $E(t)$ est dérivable pour la topologie des opérateurs et $E'(t) = B E(t)$. En effet, soit $h>0$, on utilise le fait que $\{ E(t) \ | \ t \in \mathbb{R} \}$ est un groupe abélien et que $E$ est continue en $t$
        \begin{eqnarray*}
            \vertiii{\frac{E(t+h)-E(t)}{h}-BE(t)}_1 &=& \vertiii{E(t)\Big(\frac{E(h)-I}{h}\Big)-B} \\
            &\leq& \vertiii{E(t)}_1 \vertiii{\frac{E(h)-I}{h}-B}_1 \\
            &=& \vertiii{E(t)}_1 \vertiii{BE(h)-B}_1\underset{h \to 0}{\longrightarrow} 0
        \end{eqnarray*}
        car par un changement d'indice on a obtenu
        \begin{equation*}
            \frac{E(h)-I}{h} = \frac{1}{h}\sum_{k=1}^{+\infty} \frac{h^kB^k}{k!} = \sum_{k=1}^{+\infty} \frac{h^{k-1}B^k}{k!} = \sum_{k=0}^{+\infty}\frac{h^kB^{k+1}}{k!} = BE(h)
        \end{equation*}
    \end{enumerate}
\end{remark}
Soit $u_0 \in \Hone$. On pose pour tout $t \in \bbR,\ u(t) = E(t)u_0$.
Alors par la remarque précédente, on obtient
\begin{eqnarray*}
    u'(t) = E'(t)u_0 = B E(t)u_0 = Bu(t)
\end{eqnarray*}
Le point \ref{E_groupe_abelien} de notre remarque implique aussi que $E(0) = I$, donc
\begin{eqnarray*}
    u(0) = E(0)u_0 = u_0
\end{eqnarray*}
On a bien montré l'existence d'une solution au problème \eqref{edo_ML_B}. Montrons maintenant que cette solution est unique. Soit $v(t)$ un autre solution du problème. On pose $w$ tel que
\begin{eqnarray*}
    \forall t \in \bbR ,\ w(t) := u(t) - v(t)
\end{eqnarray*}
Alors
\begin{align*}
    w'(t) &= u'(t) - v'(t) & \text{ et } &&w(0) &= u(0) -v(0)\\
          &=Bu(t) - Bv(t)  &           &&       &= u_0 - u_0\\
          &=B(u(t) - v(t)) &           &&       &= 0 \\
          &=Bw(t)          &           &&    
\end{align*}
Donc $w(t)$ est aussi solution du problème \eqref{edo_ML_B} avec pour donnée initiale $w(0)=0$. Posons 
\begin{equation*}
    \forall t \in \bbR, \ \alpha(t) := \|w(t)\|_{H^1}^2=\ps{w(t)}{w(t)}_{H^1}
\end{equation*}
Alors
\begin{eqnarray*}
    \alpha'(t) &=& \lim_{h \to 0} \frac{\alpha(t+h) - \alpha(t)}{h} \\
    &=& \limm \Big( \big\langle w(t+h),\ w(t+h)\big\rangle_{H^1} - \big\langle w(t),\ w(t)\big\rangle_{H^1} \Big) \\
    &=& \limm \Big( \big\langle w(t+h),\ w(t+h)\big\rangle_{H^1} - \big\langle w(t+h),w(t)\big\rangle_{H^1} \\
    &+& \big\langle w(t+h),w(t)\big\rangle_{H^1}-\big\langle w(t),\ w(t)\big\rangle_{H^1} \Big) \\
    &=& \limm \Big(\big\langle w(t+h), w(t+h) - w(t)\big\rangle_{H^1} + \big\langle w(t+h) - w(t), w(t)\big\rangle \Big) \\
    &=& \lim_{h \to 0} \Big(\Big\langle w(t+h), \frac{w(t+h) - w(t)}{h}\Big\rangle_{H^1} - \Big\langle \frac{w(t+h) - w(t)}{h}, w(t)\Big\rangle_{H^1} \Big) \\
    &=& 2 \big\langle w'(t), w(t) \big\rangle_{H^1} \\
    &=& 2 \big\langle Bw(t), w(t) \big\rangle_{H^1}
\end{eqnarray*}
En appliquant l'inégalité de Cauchy-Schwarz on aboutit au résultat suivant
\begin{eqnarray*}
    \abs{\alpha'(t)} &=& \abs{2 \int_{\Omega}Bw(t)w(t) + 2\int_{\Omega}D\big(Bw(t)\big) \ D\big(w(t)\big)} \\
    &\leq& 2 \norm{Bw(t)}_{L^2} \norm{w(t)}_{L^2} + 2\norm{D\big(Bw(t)\big)}_{L^2} \norm{D\big(w(t)\big)}_{L^2} \\
    &\leq& 2 \norm{B}_1 \norm{w(t)}^2_{L^2} + \norm{B}_1 \norm{D\big(w(t)\big)}^2_{L^2} \\
    &=& 2 \norm{B}_1 \norm{w(t)}^2_{H^1}\\
    &=& 2\norm{B}_1 \alpha(t)    
\end{eqnarray*}
Ce qui nous permet d'utiliser le lemme de Grönwall
\begin{eqnarray*}
    \forall t \in \bbR, \ \alpha(t) \leq \alpha(0)\exp\bigr(\abs{t}\norm{B}_1\bigl)
\end{eqnarray*}
et donc
\begin{eqnarray*}
    0 \leq \norm{w(t)}^2_{H^1} \leq \norm{w(0)}^2_{H^1} \exp\bigr(\abs{t}\norm{B}_1\bigl)
\end{eqnarray*}
Or, $w(0) = 0$, donc $\norm{w(0)}^2_{H^1}=0$, ainsi
\begin{eqnarray*}
    \forall t \in \bbR, \ \norm{w(t)}^2_1 = 0 \Longrightarrow w(t) = 0
    \Longrightarrow u(t) = v(t)
\end{eqnarray*}
D'où l'unicité de la solution.

\subsubsection*{Vérification de la solution au sens faible.}
Vérifions que pour tout $t \in \bbR,\ u(t)$ est solution du problème variationel
\begin{eqnarray*}
    \forall \varphi \in \scrD,\ \frac{\mathrm{d}}{\mathrm{d}t}b_M\big(u(t),\varphi\big) + b_L\big(u(t),\varphi\big) = 0
\end{eqnarray*}
Soit $t\in\bbR$, alors $u(t) \in \Hone$ et par densité de $\scrD$, il existe une suite $(\varphi_n)_{n\in \mathbb{N}} \subset \scrD$ tel que $\varphi_n \overunderset{\Hone}{n\rightarrow + \infty}\longrightarrow u(t)$. Par continuité, on a alors
\begin{eqnarray*}
    b_M(u'(t), \varphi) + b_L(u(t), \varphi) &=& \lim_{n \to +\infty} \big(b_M(B\varphi_n, \varphi) + b_L(\varphi_n,\varphi)\big) \\
    &=& \lim_{n\to +\infty} \big(\langle M_0B\varphi_n, \varphi\rangle_{L^2} + \langle L_0\varphi_n, \varphi\rangle_{L^2}\big) \\
    &=& \lim_{n \to +\infty} \langle M_0B\varphi_n + L_0\varphi_n,\ \varphi\rangle_{L^2} \\
    &=& 0
\end{eqnarray*}
Où pour la dernière égalité on a utilisé le fait que $M_0B = -M_0M^{-1}_0L_0 = -L_0$ et donc $M_0B\varphi_n+ L_0\varphi_n = 0$, pour tout $n \in \mathbb{N}$. L'existence et l'unicité d'une solution faible à notre problème est bien démontré.

\subsubsection{L'équation vectorielle}
Cette section fait suite à la section \ref{RS_vectorielle}. 
L'objectif ici est d'adapter le raisonnement de Showalter et Ting 
pour l'équation vectorielle. On repart donc du système 
\eqref{eq_RS}. Nous considérons un ouvert borné à la frontière 
suffisamment régulière $\Omega \subset \bbR^d$. 
Pour simplifier l'étude, on prendra des conditions de 
Dirichlet homogènes et $\bfm{f}=0$. On considère le problème 
aux limites suivant
d'inconnues $\bfm{u}$ et $\tilde{p}$
\begin{equation*}
    \left\{\begin{array}{ll}
        \partial_t\bfm{u}(t,x) - \nu\Delta \bfm{u}(t,x) - \alpha\Delta \big(\partial_t \bfm{u}(t,x)\big) = -\nabla \tilde{p}(t,x) & \forall (t,x) \in \bbR \times \Omega  \\
         \Div (\bfm{u}(t,x)) = 0 & \forall (t,x) \in \bbR \times \Omega \\
         \bfm{u}(t,x) = 0 & \forall (t,x) \in \bbR \times \partial \Omega \\
         \bfm{u}(0,x) = \bfm{u}_0(x) & \forall x \in \Omega
    \end{array}\right.
\end{equation*}
Comme dans la section précédente, on va chercher des solutions sous la forme
$\bfm{u} : \bbR \rightarrow \big(\Hone\big)^d$ et $\tilde{p} : \bbR \rightarrow L^2(\Omega)$.
\\
On considère alors le problème : trouver $(\bfm{u},\tilde{p}) \in \scrC^1\big(\bbR, \big(\Hone\big)^d\big) \times \scrC^0(\bbR, L^2(\Omega))$ tel que
\begin{equation*}
    \left\{\begin{array}{l}
        \bfm{u}'(t) - \nu\Delta \bfm{u}(t) - \alpha\Delta \big(\bfm{u}'(t)\big) = -\nabla \tilde{p}(t) \\
         \Div (\bfm{u}(t)) = 0 \\
         \bfm{u}(0) = \bfm{u}_0 
    \end{array}\right.
\end{equation*}
On introduit les opérateurs différentiels vectoriels
\begin{eqnarray*}
    \calM := I - \alpha \Delta \ \ \ \ \ &\text{ et }& \ \ \ \ \ \calL := -\nu \Delta
\end{eqnarray*}
L'objectif étant de reproduire l'analyse de la section précédente, on voudrait \og supprimer \fg{} le terme 
$\nabla \tilde{p}$. Pour ça, on introduit l'espace de Hilbert (voir Annexe \ref{AnnexeB})
\begin{equation*}
    V := \{\bfm{v} \in \big(\Hone\big)^d \ \ | \ \ \Div(\bfm{v}) = 0\}
\end{equation*}
et l'espace 
\begin{equation*}
    \calD_{\sigma}(\Omega) := \{\bfm{\varphi} \in \big(\calD(\Omega)\big)^d \ \ | \ \ \Div(\bfm{\varphi}) = 0\}
\end{equation*}

\begin{theoreme}\label{thm_densite}
    Si $\Omega \subset \bbR^d$ est un ouvert borné à la frontière suffisamment régulière alors $\calD_{\sigma}(\Omega)$ est dense dans $V$.
\end{theoreme}

Voir Annexe \ref{AnnexeB} pour la démonstration.
\\
On obtient alors un nouveau problème : trouver $(\bfm{u}, \tilde{p}) \in \scrC^1(\bbR, V) \times \scrC^0(\bbR, L^2(\Omega))$ tel que
\begin{equation}\label{pbmRSvectoriel}
    \left\{\begin{array}{l}
        \bfm{u}'(t) - \nu\Delta \bfm{u}(t) - \alpha\Delta \big(\bfm{u}'(t)\big) = -\nabla \tilde{p}(t) \\
         \bfm{u}(0) = \bfm{u}_0 
    \end{array}\right.
\end{equation}
On va établir une formulation variationnelle de \eqref{pbmRSvectoriel}. Soit $\bfm{v} \in V$.
\begin{eqnarray*}
    \int_{\Omega}\bfm{u}'(t) \cdot \bfm{v} - \nu \int_{\Omega} \Delta \bfm{u}(t) \cdot \bfm{v} - \alpha \int_{\Omega} \Delta\big(\bfm{u}'(t)\big) \cdot \bfm{v} = -\int_{\Omega} \nabla \tilde{p}(t) \cdot \bfm{v}
\end{eqnarray*}
On a en utilisant les formules de Green
\begin{eqnarray*}
    - \nu \int_{\Omega} \Delta \bfm{u}(t) \cdot \bfm{v} &=& \nu \int_{\Omega} \nabla \bfm{u}(t) : \nabla \bfm{v} -\nu \int_{\partial\Omega}\big(\nabla \bfm{u}(t) \cdot n) \cdot \bfm{v} \\
    - \alpha \int_{\Omega} \Delta\big(\bfm{u}'(t)\big) \cdot \bfm{v} &=& \alpha \int_{\Omega} \nabla \big(\bfm{u}'(t)\big) : \nabla \bfm{v} - \alpha \int_{\partial \Omega} \big(\nabla \big(\bfm{u}'(t)\big)\cdot n \big) \cdot \bfm{v}
\end{eqnarray*}
On utilise le fait que $\bfm{v} \in V$ donc en particulier $\bfm{v} \in \big(\Hone\big)^d$ donc 
\begin{eqnarray*}
    - \nu \int_{\Omega} \Delta \bfm{u}(t) \cdot \bfm{v} &=& \nu \int_{\Omega} \nabla \bfm{u}(t) : \nabla \bfm{v} \\
    - \alpha \int_{\Omega} \Delta\big(\bfm{u}'(t)\big) \cdot \bfm{v} &=& \alpha \int_{\Omega} \nabla \big(\bfm{u}'(t)\big) : \nabla \bfm{v}
\end{eqnarray*}
Enfin, on a 
\begin{equation*}
    \Div(p(t)\bfm{v}) = \nabla p(t) \cdot \bfm{v} + p(t)\Div(\bfm{v}) \Longrightarrow \nabla p(t) \cdot \bfm{v} = \Div(p(t)\bfm{v}) - p(t)\Div(\bfm{v})
\end{equation*}
d'où
\begin{equation*}
    \int_{\Omega} \nabla \tilde{p}(t) \cdot \bfm{v} = \int_{\Omega}\Div(p(t)\bfm{v}) - \int_{\Omega} p(t)\Div(\bfm{v})
\end{equation*}
et enfin, en utilisant le théorème de la divergence et le fait que $\bfm{v} \in V$ on obtient 
\begin{equation*}
    \int_{\Omega} \nabla \tilde{p}(t) \cdot \bfm{v} = \int_{\Omega}\Div(p(t)\bfm{v}) - \int_{\Omega} p(t)\Div(\bfm{v}) = \int_{\partial \Omega} p(t)(\bfm{v} \cdot n) - \int_{\Omega} p(t)\Div(\bfm{v}) = 0
\end{equation*}
On obtient alors la formulation faible : trouver $\bfm{u} \in \scrC^1\big(\bbR, V)$ tel que 
\begin{equation*}
    \forall \bfm{v} \in V, \ \int_{\Omega}\bfm{u}'(t) \cdot \bfm{v} + \nu \int_{\Omega} \nabla \bfm{u} : \nabla \bfm{v} + \alpha \int_{\Omega} \nabla \big(\bfm{u}'(t)\big) : \nabla \bfm{v} = 0
\end{equation*} 
On définit ainsi deux formes bilinéaires sur $V$ associées aux opérateurs vectoriels 
\begin{eqnarray*}
    B_{\calM}(\bfm{\varphi}, \bfm{\psi}) = \int_{\Omega} \big(\bfm{\varphi}\cdot\bfm{\psi}+\alpha\nabla \bfm{\varphi} : \nabla \bfm{\psi}\big) \ \ \ \ \ &\text{ et }& \ \ \ \ \
    B_{\calL}(\bfm{\varphi}, \bfm{\psi}) = \nu\int_{\Omega} \nabla \bfm{\varphi} : \nabla \bfm{\psi}
\end{eqnarray*}
On considère donc le problème  : trouver $\bfm{u} \in \scrC^1\big(\bbR, V)$ tel que 
\begin{equation}\label{formulation_faible_vectorielle}
    \forall t \in \bbR, \ \forall \bfm{v} \in V, \ B_{\calM}\big(\bfm{u}'(t), v\big) + B_{\calL}\big(\bfm{u}(t),v\big) = 0
\end{equation}
On cherche de nouveau à appliquer le théorème \ref{thm_lax}. On a déjà
$H=V$, il nous faut identifier qui joue le rôle de $H_0$. On voudrait reproduire le raisonnement dans le cas scalaire et ainsi identifier un espace 
de Hilbert $H_0$ tel que 
\begin{eqnarray*}
    \calD_{\sigma}(\Omega) \subset V \subset H_0 \ \ \ \ \ &\text{ et }&
    \ \ \ \ \ \overline{\calD_{\sigma}(\Omega)}^{H_0} = H_0
\end{eqnarray*}
%On choisit alors 
%\begin{equation*}
%    H_0 := \{ \bfm{f} \in \big(L^2(\Omega)\big)^d \ \ | \ \ \Div(\bfm{f}) = 0 \text{ et } \bfm{f}\cdot n = 0 \text{ sur } \partial \Omega \}
%\end{equation*}
%$H_0$ est un sous-espace fermé de $\big(L^2(\Omega)\big)^d$, en effet
%
%\begin{theoreme}
%    Si $\Omega$ (hypothèses ?) alors $\calD_{\sigma}(\Omega)$ est dense dans $H_0$.
%\end{theoreme}
%
%Voir Annexe ? pour la démonstration.
%\\
%Donc $V$ est également dense dans $H_0$.
Supposons l'existence d'un tel espace de Hilbert, dont la norme vérifie de plus l'inégalité \eqref{ineq_thm_lax}.
De manière identique que pour les formes bilinéaires $b_M$ et $b_L$, on vérifie que $B_{\calM}$ et $B_{\calL}$ sont continues et coercives sur $V$.
On peut alors appliquer le théorème \ref{thm_lax}, ainsi il existe deux opérateurs bijectifs $\calM_0 : D(\calM_0) \rightarrow H_0$ et $\calL_0 : D(\calL_0) \rightarrow H_0$ qui prolongent les opérateurs $\calM$ et $\calL$, et dont les domaines de définitions sont denses dans $V$.
On a également l'existence d'un inverse continu $\calM_0^{-1}$, on peut alors réécrire l'équation de \eqref{pbmRSvectoriel}
\begin{equation}\label{pbmRSvectoriel_operateur}
    \bfm{u}'(t) = - \calM_0^{-1}\calL_0 \bfm{u}(t)
\end{equation}
L'opérateur $-\calM_0^{-1}\calL_0 : D(\calL_0) \rightarrow D(\calM_0)$ est une bijection et, de façon similaire à l'opérateur $-M_0^{-1}L_0$, on montre que $-\calM_0^{-1}\calL_0$ est continu sur $V$ et ainsi on applique le théorème \ref{thm_prolongement} et on prolonge de manière unique $-\calM_0^{-1}\calL_0$ en l'opérateur $\calB : V \rightarrow V$.
Cet opérateur $\calB$ est inversible d'inverse continu sur $V$.
Ainsi, \eqref{pbmRSvectoriel_operateur} devient 
\begin{equation}\label{pbmRSvectoriel_opB}
    \bfm{u}'(t) = \calB \bfm{u}(t)
\end{equation}
On veut maintenant définir $\exp(t\calB)$. Comme pour l'équation 
scalaire $\calB \in \mathscr{L}(V)$. Et $\mathscr{L}(V)$ est 
un espace vectoriel complet, car $V$ est complet, on peut donc 
définir l'exponentielle de la façon suivante
\begin{eqnarray*}
\forall t \in \mathbb{R},\  \calE(t) := \exp(t\calB) = \sum_{k=0}^{+\infty} \frac{(t\calB)^k}{k!}
\end{eqnarray*}
À l'aide de la norme opérateur l'on obtient ainsi
\begin{equation*}
    \forall t\in \mathbb{R}, \ \vertiii{\calE(t)}_1 \leq \exp\big(\abs{t}\vertiii{\calB}_1\big)
\end{equation*}
Donc $\{ \calE(t) \ | \ t \in \mathbb{R} \} $ est un ensemble d'opérateurs continus de $V$.

\begin{remark}
    La remarque \ref{remarque_operateur} s'applique à $\calE$.
\end{remark}

Soit $\bfm{u}_0 \in V$. On pose pour tout $t \in \bbR,\ \bfm{u}(t) = \calE(t)\bfm{u}_0$.
Alors par la remarque précédente, on obtient
\begin{eqnarray*}
    \bfm{u}'(t) = \calE'(t)\bfm{u}_0 = \calB \calE(t)\bfm{u}_0 = \calB \bfm{u}(t)
\end{eqnarray*}
Notre remarque précédente implique aussi que $\calE(0) = I$, donc
\begin{eqnarray*}
    \bfm{u}(0) = \calE(0)\bfm{u}_0 = \bfm{u}_0
\end{eqnarray*}
d'où l'existence. Pour montrer l'unicité, il suffit de reproduire 
le raisonnement scalaire. Soit $\bfm{v}(t)$ un autre solution du 
problème \eqref{pbmRSvectoriel_opB}. On pose $\bfm{w}$ tel que
\begin{eqnarray*}
    \forall t \in \bbR ,\ \bfm{w}(t) := \bfm{u}(t) - \bfm{v}(t)
\end{eqnarray*}
Alors $\bfm{w}(t)$ est aussi solution du problème \eqref{pbmRSvectoriel_opB} avec pour donnée initiale $\bfm{w}(0)=0$. Posons 
\begin{equation*}
    \forall t \in \bbR, \ \alpha(t) := \|\bfm{w}(t)\|_{H^1}^2=\ps{\bfm{w}(t)}{\bfm{w}(t)}_{H^1}
\end{equation*}
Alors
\begin{eqnarray*}
    \alpha'(t)= 2 \big\langle \calB \bfm{w}(t), \bfm{w}(t) \big\rangle_{H^1}
\end{eqnarray*}
En appliquant l'inégalité de Cauchy-Schwarz on aboutit au résultat suivant
\begin{eqnarray*}
    \abs{\alpha'(t)} \leq 2\norm{\calB}_1 \alpha(t)    
\end{eqnarray*}
Ce qui nous permet d'utiliser le lemme de Grönwall
\begin{eqnarray*}
    \forall t \in \bbR, \ \alpha(t) \leq \alpha(0)\exp\big(\abs{t}\norm{\calB}_1\big)
\end{eqnarray*}
et donc
\begin{eqnarray*}
    0 \leq \norm{\bfm{w}(t)}^2_{H^1} \leq \norm{\bfm{w}(0)}^2_{H^1} \exp\big(\abs{t}\norm{\calB}_1\big)
\end{eqnarray*}
Or, $\bfm{w}(0) = 0$, donc $\norm{\bfm{w}(0)}^2_{H^1}=0$, ainsi
\begin{eqnarray*}
    \forall t \in \bbR, \ \norm{\bfm{w}(t)}^2_1 = 0 \Longrightarrow \bfm{w}(t) = 0
    \Longrightarrow \bfm{u}(t) = \bfm{v}(t)
\end{eqnarray*}
D'où l'unicité de la solution. Vérifions que pour tout $t \in \bbR,\ \bfm{u}(t)$ est solution du problème variationnel
\begin{eqnarray*}
    \forall \bfm{\varphi} \in \calD_{\sigma}(\Omega),\ \frac{\mathrm{d}}{\mathrm{d}t}B_{\calM}\big(\bfm{u}(t),\bfm{\varphi}\big) + B_{\calL}\big(\bfm{u}(t),\bfm{\varphi}\big) = 0
\end{eqnarray*}
Soit $t\in\bbR$, alors $\bfm{u}(t) \in V$ et par densité de $\calD_{\sigma}(\Omega)$, il existe une suite $(\bfm{\varphi}_n)_{n\in \mathbb{N}} \subset \calD_{\sigma}(\Omega)$ tel que $\bfm{\varphi}_n \overunderset{V}{n\rightarrow + \infty}\longrightarrow \bfm{u}(t)$. Par continuité, on a alors
\begin{eqnarray*}
    B_{\calM}(\bfm{u}'(t), \bfm{\varphi}) + B_{\calL}(\bfm{u}(t), \bfm{\varphi}) = 0
\end{eqnarray*}
L'existence et l'unicité d'une solution faible à notre problème 
est bien démontré.
\paragraph{Détermination de la pression.} Il faut maintenant revenir au problème initial et déterminer $\tilde{p}$. Pour ça, on se base sur le raisonnement mené dans la section $2$ de \cite{Temam1977}. Pour tout $t \in \bbR$, $\bfm{u}(t) \in V$ est solution de \eqref{formulation_faible_vectorielle} donc pour tout $t \in \bbR$, $\bfm{u}(t)$ est solution au sens des distributions. On a donc pour tout $t \in \bbR$
\begin{equation*}
    \forall \bfm{v} \in \calD_{\sigma}(\Omega), \ \langle \bfm{u}'(t) - \nu \Delta \bfm{u}(t) - \alpha \Delta \bfm{u}'(t), \bfm{v} \rangle = 0
\end{equation*}
On utilise 
\begin{proposition}[Lemme de De Rham]\label{lemme_De_Rham}
    Soient un ouvert $\Omega \subset \bbR^d$ et $\bfm{f} \in \big(\calD'(\Omega)\big)^d$. Alors il existe 
    $p \in \calD'(\Omega)$ tel que $\bfm{f} = \nabla p$ si et seulement si 
    \begin{equation*}
        \forall \bfm{v} \in \calD_{\sigma}(\Omega), \ \langle \bfm{f}, \bfm{v} \rangle =0
    \end{equation*}
\end{proposition}
On applique ce résultat pour $\bfm{u}'(t) - \nu \Delta \bfm{u}(t) - \alpha \Delta\bfm{u}'(t) \in V'$, ce qui nous donne l'existence de $p \in \calD'(\Omega)$ tel que 
\begin{equation*}
    \bfm{u}'(t) - \nu \Delta \bfm{u}(t) - \alpha\Delta \bfm{u}'(t) = -\nabla p
\end{equation*}
au sens des distributions. Et ceci étant vrai pour tout $t \in \bbR$, on définit une fonction $\tilde{p}:\bbR \rightarrow \calD'(\Omega)$ telle que 
\begin{equation*}
    \bfm{u}'(t) - \nu \Delta \bfm{u}(t) - \alpha\Delta \bfm{u}'(t) = -\nabla \tilde{p}(t)
\end{equation*}

\begin{remark}
    L'idée générale du raisonnement pour montrer l'existence et l'unicité d'une solution $\bfm{u}$ et l'existence d'une solution $\tilde{p}$ à l'équation vectorielle \eqref{RS_vectorielle} est ici. Il reste cependant de nombreux points à détailler
    \begin{enumerate}
        \item il nous faudrait identifier l'espace de Hilbert $H_0$, il s'agit certainement d'un sous-ensemble de $\big(L^2(\Omega)\big)^d$ à divergence nulle mais avec une condition supplémentaire, car la seule condition de divergence nulle ne permet pas de conclure ;
        \item la démonstration du théorème \ref{thm_densite} est incomplète et ne nous permet pas de conclure quant à la densité de $\calD_{\sigma}(\Omega)$ dans $V$ ;
        \item une démonstration du lemme de De Rham \ref{lemme_De_Rham} qui est un résultat que nous admettons ici ;
        \item justifier l'inclusion $V' \subset \big(\calD'(\Omega)\big)^d$ afin d'utiliser le lemme de De Rham ;
        \item montrer que pour tout $t \in \bbR$, $\tilde{p}(t) \in L^2(\Omega)$, la proposition 1.2 de \cite{Temam1977} permettrait d'établir cela.
    \end{enumerate}
\end{remark}


\newpage